\section{System boundary and comparison basis}
The final candidate boundary follows the source INL process:
\[
\text{natural gas + steam}\rightarrow\text{SMR}\rightarrow\text{WGS}
\rightarrow\text{CO$_2$ capture}\rightarrow\text{PSA}\rightarrow\text{H$_2$}.
\]
High-temperature nuclear heat replaces reformer firing in the final Case-6 architecture. The lifecycle boundary adds upstream natural-gas supply, nuclear heat allocation, process electricity and CO$_2$ transport/storage. The economic boundary adds natural gas, full reactor cogeneration burden, incremental CCS capital, integration/site allowance and T\&S, with net exported electricity valued separately.

The principal baseline is INL's conventional no-capture SMR at the same 130 MMSCFD H$_2$ output~\cite{inltev953}. This provides a source-consistent direct-emissions and natural-gas comparison rather than scaling an unrelated plant. IEAGHG Case 1A remains an external conventional-SMR+CCS comparator, but its Singapore economics are not used to overwrite the final INL common-output baseline.
